
\documentclass{article}
\usepackage{colortbl}
\usepackage{makecell}
\usepackage{multirow}
\usepackage{supertabular}

\begin{document}

\newcounter{utterance}

\centering \large Interaction Transcript for game `dond', experiment `semi\_en', episode 2 with reflexion{-}Qwen3.5{-}9B.
\vspace{24pt}

{ \footnotesize  \setcounter{utterance}{1}
\setlength{\tabcolsep}{0pt}
\begin{supertabular}{c@{$\;$}|p{.15\linewidth}@{}p{.15\linewidth}p{.15\linewidth}p{.15\linewidth}p{.15\linewidth}p{.15\linewidth}}
    \# & $\;$A & \multicolumn{4}{c}{Game Master} & $\;\:$B\\
    \hline

    \theutterance \stepcounter{utterance}  
    & & \multicolumn{4}{p{0.6\linewidth}}{
        \cellcolor[rgb]{0.9,0.9,0.9}{
            \makecell[{{p{\linewidth}}}]{
                \texttt{\tiny{[P1$\langle$GM]}}
                \texttt{You have attempted this same task before and did not succeed. These are the plans you wrote after those attempts.} \\
\texttt{Trial \#1: New plan: I failed because I misunderstood the timing of the "secret proposal" mechanism. I treated the initial negotiation messages as binding offers that would be compared directly, rather than realizing that the secret proposal is a separate, final action that overrides or complements the public negotiation. By sending "Stick 1, Pepper 1" publicly, I created a conflict with the other player's identical public message, resulting in zero points. Next time, I will use the initial messages solely to gather information about the other player's preferences and to signal my own high{-}value demands without committing to a specific bundle. I will save the actual allocation for the secret proposal phase, ensuring my secret proposal is perfectly complementary to the other player's likely public stance (e.g., if they ask for the stick, I will secretly propose the carrot and peppers). Specifically, I will start by stating my valuation strategy, then wait for their response, and finally submit a secret proposal like "{[}Proposal: 1 carrot, 2 peppers{]}" only after confirming the other player is not claiming those specific items.} \\
\\ 
\texttt{Trial \#2: New plan: I failed because I publicly claimed the stick in my first message ("I want the stick"), which directly conflicted with the other player's identical public claim, resulting in a zero{-}score outcome. The rules state that public messages are transmitted directly and can create conflicts, whereas the secret proposal is the only mechanism to guarantee a non{-}conflicting allocation. My strategy of using public messages to "signal interest" without binding commitment was flawed because my specific phrasing ("I want the stick") was interpreted as a binding claim. Next time, I will strictly avoid mentioning specific item allocations in public messages. Instead, I will only communicate my valuation hierarchy (e.g., "Stick is most valuable to me") to gauge the opponent's stance. I will then immediately submit a secret proposal that claims the high{-}value items I need (the stick) while explicitly excluding items the opponent might claim, ensuring our proposals are complementary. Specifically, I will send a public message stating my values, wait for their response, and then submit: "{[}Proposal: 1 stick, 0 carrot, 0 peppers{]}" (or similar) only after confirming they are not claiming the stick, or simply claim the stick in the secret proposal if they claim it publicly, knowing the secret proposal overrides the public conflict.} \\
\\ 
\\ 
\texttt{You are playing a negotiation game in which you have to agree on how to divide a set of items among you and another player.} \\
\\ 
\texttt{Rules:} \\
\texttt{(a) You and the other player are given a set of items. Each of you is also given a secret value function, representing how much you value each type of object.} \\
\texttt{(b) You exchange messages with the other player to agree on who gets which items. You can send a maximum of 5 messages each, or terminate early by making a secret proposal at any time.} \\
\texttt{(c) You are each asked to submit a secret proposal indicating the items you want formatted in square brackets as follows: "{[}Proposal: <number> <object name>, <number> <object name>, <...>{]}"} \\
\texttt{(d) If your proposals are complementary, i.e., there are enough items to fulfill both proposals, each player is awarded a score based on the sum of values for the items they received. Otherwise, both of you get zero points.} \\
\texttt{(e) Your goal is to maximize the score you receive.} \\
\\ 
\texttt{Let us start.} \\
\\ 
\texttt{The set of available items is:} \\
\\ 
\texttt{1 carrot, 1 stick, 3 peppers} \\
\\ 
\texttt{Your secret value function is:} \\
\\ 
\texttt{carrot: 1, stick: 3, pepper: 2} \\
\\ 
\texttt{IMPORTANT: Your messages, unless it is the secret proposal, are directly transmitted to the other player, so do not include any response to the rules or text announcing your message. To make a secret proposal, use the indicated format. Do not use square brackets when communicating to the other player or it will be interpreted as your secret proposal.} \\
            }
        }
    }
    & \\ \\

    \theutterance \stepcounter{utterance}  
    & \multicolumn{4}{p{0.6\linewidth}}{
        \cellcolor[rgb]{0.9,0.9,0.9}{
            \makecell[{{p{\linewidth}}}]{
                \texttt{\tiny{[P1$\rangle$GM]}}
                \texttt{Stick is most valuable to me, followed by peppers, then carrots. I am interested in a deal where I get the stick and you get the peppers, or we split the peppers. Let me know your preferences.} \\
            }
        }
    }
    & & \\ \\

    \theutterance \stepcounter{utterance}  
    & & & \multicolumn{2}{p{0.3\linewidth}}{
        \cellcolor[rgb]{0.9,0.9,0.9}{
            \makecell[{{p{\linewidth}}}]{
                \texttt{\tiny{[GM$|$GM]}}
                \texttt{continue} \\
            }
        }
    }
    & & \\ \\

    \theutterance \stepcounter{utterance}  
    & & & \multicolumn{2}{p{0.3\linewidth}}{
        \cellcolor[rgb]{0.9,0.9,0.9}{
            \makecell[{{p{\linewidth}}}]{
                \texttt{\tiny{[GM$|$GM]}}
                \texttt{Stick is most valuable to me, followed by peppers, then carrots. I am interested in a deal where I get the stick and you get the peppers, or we split the peppers. Let me know your preferences.} \\
            }
        }
    }
    & & \\ \\

    \theutterance \stepcounter{utterance}  
    & & \multicolumn{4}{p{0.6\linewidth}}{
        \cellcolor[rgb]{0.9,0.9,0.9}{
            \makecell[{{p{\linewidth}}}]{
                \texttt{\tiny{[GM$\rangle$P2]}}
                \texttt{You have attempted this same task before and did not succeed. These are the plans you wrote after those attempts.} \\
\texttt{Trial \#1: New plan: I failed because I misunderstood the critical rule regarding message formatting. I sent "Stick 1, Pepper 1" as a public message, but the rules explicitly state that only the specific format "{[}Proposal: ...{]}" triggers a secret proposal; otherwise, messages are directly transmitted to the opponent. By not using the secret proposal format immediately, I revealed my high valuation of the stick to the opponent, allowing them to counter with a proposal that left me with only low{-}value items (peppers worth 0 to me). Next time, I will immediately use the secret proposal format "{[}Proposal: 1 stick, 1 pepper{]}" in my first turn to hide my true valuations and force the opponent to make a blind guess, maximizing my chance of securing the high{-}value stick without revealing my strategy.} \\
\\ 
\texttt{Trial \#2: New plan: I failed because I revealed my high valuation of the stick in a public message ("I want the stick"), allowing the opponent to anticipate my move and likely secure it in their secret proposal. Since the stick is my highest value item (3 points) and the opponent values it most (3 points), leaving it unsecured in a public bid was a strategic error. Next time, I will immediately submit my secret proposal in the exact required format "{[}Proposal: 1 stick, 1 pepper{]}" on the very first turn. This hides my true preferences, prevents the opponent from tailoring their counter{-}proposal against me, and maximizes the probability that our proposals are complementary, ensuring I secure the stick and the pepper for a total score of 2 (stick: 3 + pepper: 0 = 3, assuming the opponent takes the carrot).} \\
\\ 
\\ 
\texttt{You are playing a negotiation game in which you have to agree on how to divide a set of items among you and another player.} \\
\\ 
\texttt{Rules:} \\
\texttt{(a) You and the other player are given a set of items. Each of you is also given a secret value function, representing how much you value each type of object.} \\
\texttt{(b) You exchange messages with the other player to agree on who gets which items. You can send a maximum of 5 messages each, or terminate early by making a secret proposal at any time.} \\
\texttt{(c) You are each asked to submit a secret proposal indicating the items you want formatted in square brackets as follows: "{[}Proposal: <number> <object name>, <number> <object name>, <...>{]}"} \\
\texttt{(d) If your proposals are complementary, i.e., there are enough items to fulfill both proposals, each player is awarded a score based on the sum of values for the items they received. Otherwise, both of you get zero points.} \\
\texttt{(e) Your goal is to maximize the score you receive.} \\
\\ 
\texttt{Let us start.} \\
\\ 
\texttt{The set of available items is:} \\
\\ 
\texttt{1 carrot, 1 stick, 3 peppers} \\
\\ 
\texttt{Your secret value function is:} \\
\\ 
\texttt{carrot: 8, stick: 2, pepper: 0} \\
\\ 
\texttt{IMPORTANT: Your messages, unless it is the secret proposal, are directly transmitted to the other player, so do not include any response to the rules or text announcing your message. To make a secret proposal, use the indicated format. Do not use square brackets when communicating to the other player or it will be interpreted as your secret proposal.} \\
\\ 
\\ 
\texttt{Stick is most valuable to me, followed by peppers, then carrots. I am interested in a deal where I get the stick and you get the peppers, or we split the peppers. Let me know your preferences.} \\
            }
        }
    }
    & \\ \\

    \theutterance \stepcounter{utterance}  
    & & & \multicolumn{4}{p{0.6\linewidth}}{
        \cellcolor[rgb]{0.9,0.9,0.9}{
            \makecell[{{p{\linewidth}}}]{
                \texttt{\tiny{[GM$\langle$P2]}}
                \texttt{I am interested in securing the stick for myself, as it is my highest priority. I value the stick at 2 points, while the carrot is worth 8 points to me (though I may not need it if I get the stick and peppers). The peppers are worth 0 to me, so I am indifferent to who gets them, provided I secure the stick.} \\
\\ 
\texttt{To maximize our chances of a successful agreement where I get the stick and you get the peppers (or we split the peppers), I propose the following division:} \\
\\ 
\texttt{I take the stick and the carrot.} \\
\texttt{You take the 3 peppers.} \\
\\ 
\texttt{This way, I secure my most valuable item (the stick), and you receive the items you likely value more (the peppers). Let's see if this works for you.} \\
\\ 
\texttt{{[}Proposal: 1 stick, 1 carrot{]}} \\
            }
        }
    }
     \\ \\

    \theutterance \stepcounter{utterance}  
    & & & \multicolumn{2}{p{0.3\linewidth}}{
        \cellcolor[rgb]{0.9,0.9,0.9}{
            \makecell[{{p{\linewidth}}}]{
                \texttt{\tiny{[GM$|$GM]}}
                \texttt{proposal} \\
            }
        }
    }
    & & \\ \\

    \theutterance \stepcounter{utterance}  
    & & & \multicolumn{2}{p{0.3\linewidth}}{
        \cellcolor[rgb]{0.9,0.9,0.9}{
            \makecell[{{p{\linewidth}}}]{
                \texttt{\tiny{[GM$|$GM]}}
                [1, 1, 0]
            }
        }
    }
    & & \\ \\

    \theutterance \stepcounter{utterance}  
    & & \multicolumn{4}{p{0.6\linewidth}}{
        \cellcolor[rgb]{0.9,0.9,0.9}{
            \makecell[{{p{\linewidth}}}]{
                \texttt{\tiny{[P1$\langle$GM]}}
                \texttt{The other player has submitted a secret proposal. It is now your turn to submit a proposal of your own.} \\
            }
        }
    }
    & \\ \\

    \theutterance \stepcounter{utterance}  
    & \multicolumn{4}{p{0.6\linewidth}}{
        \cellcolor[rgb]{0.9,0.9,0.9}{
            \makecell[{{p{\linewidth}}}]{
                \texttt{\tiny{[P1$\rangle$GM]}}
                \texttt{{[}Proposal: 1 stick, 0 carrot, 2 peppers{]}} \\
            }
        }
    }
    & & \\ \\

    \theutterance \stepcounter{utterance}  
    & & & \multicolumn{2}{p{0.3\linewidth}}{
        \cellcolor[rgb]{0.9,0.9,0.9}{
            \makecell[{{p{\linewidth}}}]{
                \texttt{\tiny{[GM$|$GM]}}
                \texttt{proposal} \\
            }
        }
    }
    & & \\ \\

    \theutterance \stepcounter{utterance}  
    & & & \multicolumn{2}{p{0.3\linewidth}}{
        \cellcolor[rgb]{0.9,0.9,0.9}{
            \makecell[{{p{\linewidth}}}]{
                \texttt{\tiny{[GM$|$GM]}}
                [0, 1, 2]
            }
        }
    }
    & & \\ \\

    \theutterance \stepcounter{utterance}  
    & & & \multicolumn{2}{p{0.3\linewidth}}{
        \cellcolor[rgb]{0.9,0.9,0.9}{
            \makecell[{{p{\linewidth}}}]{
                \texttt{\tiny{[GM$|$GM]}}
                [[0, 1, 2], [1, 1, 0]]
            }
        }
    }
    & & \\ \\

    \theutterance \stepcounter{utterance}  
    & \multicolumn{4}{p{0.6\linewidth}}{
        \cellcolor[rgb]{0.95,0.95,0.95}{
            \makecell[{{p{\linewidth}}}]{
                \texttt{\tiny{[P2$\rightarrow$P1]}}
                \texttt{Post{-}episode Reflexion (not sent to GM):} \\
\texttt{New plan: I failed because I submitted a public message containing a specific allocation ("I am interested in a deal where I get the stick..."), which the game interpreted as a binding claim that conflicted with the opponent's secret proposal, resulting in a zero score. The rules explicitly state that public messages are transmitted directly and can cause conflicts, while only the secret proposal format guarantees a valid outcome. My strategy of using public chat to negotiate specific splits was flawed. Next time, I will strictly adhere to the rule: public messages must contain **only** my valuation hierarchy (e.g., "Stick is most valuable, then peppers, then carrots") with **no** mention of specific item counts or allocation requests. I will use the public channel solely to signal my priorities and wait for the opponent's response. Once the opponent has made their move, I will immediately submit my secret proposal using the exact format "{[}Proposal: <items>{]}" to claim the high{-}value stick and the peppers, ensuring my proposal is complementary to whatever they publicly stated or secretly proposed, thereby securing the maximum possible score without triggering a conflict.} \\
            }
        }
    }
    & & \\ \\

    \theutterance \stepcounter{utterance}  
    & & & \multicolumn{4}{p{0.6\linewidth}}{
        \cellcolor[rgb]{0.95,0.95,0.95}{
            \makecell[{{p{\linewidth}}}]{
                \texttt{\tiny{[P1$\rightarrow$P2]}}
                \texttt{Post{-}episode Reflexion (not sent to GM):} \\
\texttt{New plan: I failed because I revealed my true valuation of the stick (2 points) and my indifference to peppers in a public message, allowing the opponent to anticipate my needs and likely secure the stick in their secret proposal. Since the stick is my highest{-}value item, exposing this preference before making a secret offer was a critical strategic error. Next time, I will immediately submit my secret proposal in the exact required format "{[}Proposal: 1 stick, 1 carrot{]}" on the very first turn without any preceding public discussion. This hides my true preferences, prevents the opponent from tailoring a counter{-}proposal to take the stick, and maximizes the probability that our proposals are complementary, ensuring I secure the stick and the carrot for a total score of 10.} \\
            }
        }
    }
     \\ \\

\end{supertabular}
}

\end{document}
